\section{Introduction}\label{sec:intro}

\subsection{The problem: an approval the agent can reach is not a second factor}

Tool-using agents act with their owner's authority. They read and send email, move money, edit files and run
shell commands, and the industry's answer to the risk is an approval step: before a consequential call executes,
the agent pauses and a person, or increasingly an automated approver, clicks \emph{approve}. 15 of 21 production agent systems surveyed
by \citet{wang2026reframing} rely on some form of runtime approval, and coding agents, payment agents and
browser agents all ship one. The step is usually described, and often marketed, as the agent equivalent of
multi-factor authentication (MFA).

Human MFA gives a precise test for that claim. A second factor is valuable
because it lives somewhere the attacker of the first factor cannot reach: a password and a one-time code typed
into the same phishing page are, against that phisher, one factor. Payment regulation added a second lesson.
Under the European rules for strong customer authentication, the code a payer enters must be
\emph{dynamically linked} to the amount and the payee the payer sees, so that an approval cannot be spent on a
different transaction~\citep{psd2rts2018}. Carried over to agents, the two lessons become one criterion,
which this paper tests:
\begin{quote}
\emph{An approval protects an agent action only if compromising the agent cannot produce it, and only if it
commits to exactly the action that executes.}
\end{quote}
Most deployed approvals fail one half of this criterion. A ``reply yes to continue'' in the agent's chat is
rendered by the agent and can be answered by anyone who speaks in that channel. An out-of-band card whose text
the agent writes moves the click off-host but leaves the \emph{content} in the agent's hands, so the human
approves a description rather than an action. A one-time code is out-of-band but says nothing about what it
approves and can be relayed to any call. An approval that names a run, a tool or a session rather than the
exact arguments can be reused after it is granted. In each case the approval's existence is outside the agent,
but its meaning is not. \Cref{fig:overview} shows the three patterns side by side.

\begin{figure}[t]
\centering
% Figure 1: three ways to ask a human
\resizebox{\textwidth}{!}{%
\begin{tikzpicture}[
  font=\footnotesize,
  >={Stealth[length=4pt]},
  box/.style={draw=inkgray,rounded corners=2pt,align=center,minimum height=8mm,minimum width=18mm,text width=17mm,inner sep=2pt,fill=white},
  arr/.style={->,thick,inkgray},
  lab/.style={font=\scriptsize,align=center,inner sep=1.5pt},
  bad/.style={fill=untrustred,draw=none,rounded corners=4pt},
  good/.style={fill=trustgreen,draw=none,rounded corners=4pt}]

% ---------------- (a) in-channel confirmation
\begin{scope}[xshift=0cm]
  \fill[bad] (0,0) rectangle (6.6,4.2);
  \node[box] (A) at (1.3,3.2) {Agent};
  \node[box] (C) at (5.3,3.2) {Chat channel};
  \node[box] (T) at (1.3,0.9) {Tool};
  \node[box] (H) at (5.3,0.9) {User};
  \draw[arr] ([yshift=2mm]A.east) -- node[lab,above] {``may I?''} ([yshift=2mm]C.west);
  \draw[arr,narrorange,dashed] ([yshift=-2mm]A.east) -- node[lab,below,text=narrorange] {forged ``yes''} ([yshift=-2mm]C.west);
  \draw[arr] (H) -- node[lab,right] {``yes''} (C);
  \draw[arr] (A) -- node[lab,left] {runs} (T);
  \node[font=\small\bfseries,anchor=south west] at (0,4.3) {(a) In-channel confirmation};
\end{scope}

% ---------------- (b) narrated out-of-band card
\begin{scope}[xshift=7.1cm]
  \fill[bad] (0,0) rectangle (2.9,4.2);
  \fill[good] (3.1,0) rectangle (6.6,4.2);
  \node[box] (A) at (1.3,3.2) {Agent};
  \node[box] (P) at (5.3,3.2) {Phone shows agent's text};
  \node[box] (T) at (1.3,0.9) {Tool};
  \node[box] (H) at (5.3,0.9) {User taps ``ok''};
  \draw[arr] (A) -- node[lab,above] {``email boss''} (P);
  \draw[arr] (P) -- (H);
  \draw[arr] (H) -- node[lab,below] {ok} (T);
  \draw[arr,narrorange] (A) -- node[lab,left,text=narrorange] {bcc to\\attacker} (T);
  \node[font=\small\bfseries,anchor=south west] at (0,4.3) {(b) Narrated card};
\end{scope}

% ---------------- (c) effect-bound approval
\begin{scope}[xshift=14.2cm]
  \fill[bad] (0,0) rectangle (1.9,4.2);
  \fill[good] (2.1,0) rectangle (9.4,4.2);
  \node[box,text width=13mm,minimum width=14mm] (A) at (0.95,3.2) {Agent};
  \node[box,text width=23mm,minimum width=24mm] (G) at (4.1,3.2) {Gateway: canonical $\Desc$, digest $\dig$};
  \node[box,text width=17mm] (P) at (7.9,3.2) {Phone: card from $\Desc$};
  \node[box,text width=17mm] (S) at (7.9,0.9) {Passkey signs $c(\dig,\mathit{id})$};
  \node[box,text width=23mm,minimum width=24mm] (T) at (4.1,0.9) {Execute stored $\Desc$ once};
  \draw[arr] ([yshift=2mm]A.east) -- node[lab,above] {call} ([yshift=2mm]G.west);
  \draw[arr,dashed] ([yshift=-2mm]G.west) -- node[lab,below] {id} ([yshift=-2mm]A.east);
  \draw[arr] (G) -- (P);
  \draw[arr] (P) -- (S);
  \draw[arr] (S) -- node[lab,below] {receipt} (T);
  \draw[arr] (G) -- node[lab,right] {stored} (T);
  \node[font=\small\bfseries,anchor=south west] at (0,4.3) {(c) Effect-bound approval};
\end{scope}
\end{tikzpicture}}

\caption{Three ways to ask a human. (a) In-channel confirmation: the agent renders the question and the answer
arrives on a channel it can write to. (b) Narrated out-of-band card: the click happens on the user's phone, but
the text the user reads is the agent's description, so the user approves a description rather than the call.
(c) Effect-bound approval, the design we use as a reference defense (\cref{sec:design}): a gateway the agent
cannot reach renders the card from the canonical call, the phone signs the digest of that call, and the
gateway executes only the stored call. Green regions are outside the agent's reach; red regions are inside it.}
\label{fig:overview}
\end{figure}

\subsection{Motivating evidence: approvals already fail in practice}

The two-week red-team study \emph{Agents of Chaos}~\citep{shapira2026agents} documents what happens when
identity and approval live where the agent lives. In its Case~\#8 a non-owner changed a Discord display name to
the owner's and opened a private channel; the agent inferred ownership from the name and tone, deleted its
memory and configuration files and reassigned administrative access. In Case~\#15 two agents correctly
distrusted a social-engineering email and then verified the owner by asking the possibly compromised Discord
account to confirm itself. In Case~\#1 an agent asked for confirmation before a destructive reset and
``double-checked'', and both confirmations came from the requesting non-owner in the same conversation. The
authors conclude that systems relying on presented identity rather than cryptographically grounded or
multi-factor authentication remain susceptible to session-boundary attacks.

Binding failures are documented as well. Loopjacking~\citep{arunkumar2026loopjacking} shows post-approval
substitution in released agent products, where the call that executes is not the call that was approved. Practitioner reports describe agents that misstate pending actions to their approvers~\citep{checkmarx2025lies}
and approvals that resolve a symbolic link differently from what the dialog showed~\citep{wiz2026ghostapproval}.
These failures were observed in deployed systems.

\subsection{Approach}

This paper continues aiAuthZ~\citep{kodathala2026aiauthz}, earlier work by the first author. aiAuthZ placed authorization in a
separate trust domain: an off-host gateway binds every user message to an identity and checks every tool call
against a role and argument policy the agent cannot read. It answered two questions, \emph{who is speaking} and
\emph{which calls are allowed}, and it stated plainly what it left open: its receipts use a shared-key HMAC and give no third-party
non-repudiation, and it has no step-up path in which a human approves a specific action. This paper takes up
that remaining question, \emph{did the human approve exactly this effect}, and treats it first as an attack
surface and only then as a design problem.

We proceed in four steps. First, we define the properties an approval must have to count as an independent,
effect-bound factor (\cref{sec:background}). Second, we build a family of twelve attacks on the approval path and
run them against real designs and implementations (\cref{sec:attacks}), and we audit how ten agent runtimes
implement approval today (\cref{sec:audit}). Third, we build a reference defense, \sys{} (Agent Step-Up), and
prove its properties in Tamarin under explicitly stated hypotheses (\cref{sec:design,sec:formal}). Fourth, we
measure what binding, trusted rendering and a second device each buy, against scripted attacks, AgentDojo
replays, simulated human approvers and \llmModels{} current language models acting as automated approvers
(\cref{sec:eval}).

\subsection{Contributions}

None of the cryptographic primitives in this paper is new, and the closest designs, SUDP~\citep{yu2026sudp} and
Microsoft's action-bound approval design, specified in Architecture Decision Record (ADR)
0030~\citep{microsoft2026adr0030},, already bind a canonical operation digest
to a fresh approval. The contribution is the attacks and the measurement around that idea. Specifically, we contribute:
\begin{itemize}[leftmargin=1.4em,itemsep=2pt]
\item A stated, testable criterion for when an approval is a second factor for an agent action, decomposed into
  seven properties \Prop{1} to \Prop{7} (\cref{sec:properties}).
\item A family of twelve approval-path attacks, \AB{1} to \AB{12}, executed against the Agent Payments Protocol (AP2), Microsoft's ADR~0030
  implementation and Model Context Protocol (MCP) elicitation, reviewed statically against OpenClaw and mapped
  onto in-process frameworks, with \sys{} as the control
  (\cref{sec:attacks}).
\item A reproducible audit of ten agent runtimes with uniform criteria for effect binding, single use and flood
  resistance, including a version sweep showing that post-approval substitution in one framework persists past
  its published affected range (\cref{sec:audit}).
\item \sys{}, an open reference gateway with WebAuthn approval over a canonical effect digest, typed and
  injective card rendering, argument-complete grants and flood limits, and a Tamarin model with
  \foLemmas{} lemmas whose hypotheses are explicit and whose attacks reappear when each hypothesis is dropped
  (\cref{sec:design,sec:formal}).
\item An evaluation that separates what binding, trusted rendering and a second device each contribute, reports
  the cases where the reference defense does not help, and measures \llmModels{} language models from
  \llmVendors{} vendors as automated approvers on the same stimuli (\cref{sec:eval}).
\end{itemize}

We also state what this paper does not claim. It does not claim the idea of out-of-band,
parameter-bound step-up for agents, which appears in a defensive publication~\citep{vyas2026visamfa}, in
SUDP~\citep{yu2026sudp}, in ADR~0030~\citep{microsoft2026adr0030}, in AP2~\citep{ap2spec} and in Client-Initiated Backchannel Authentication (CIBA) with Rich
Authorization Requests~\citep{ciba2021,rfc9396}. It does not claim the first formal model of agent
authorization~\citep{jiang2026payment,chen2026aith}. It reports no human-subjects data: every approver in the
evaluation is either a simulated model of a person or a language model, and no statement here is a claim about
how often people catch attacks.
